\documentclass[11pt]{article}
\usepackage{hyperref}
\setlength{\oddsidemargin}{-0.5in}
\setlength{\evensidemargin}{\oddsidemargin}
\setlength{\textwidth}{7.0in}
\setlength{\textheight}{8.3in}
\setlength{\topmargin}{-0.7in}
%\input{macros.tex}
\begin{document}

{\Large COSC 417 \hspace{2cm} Assignment: Programming Task 1 }
\vspace{0.5cm}

\emph{Instructions:}

(a) Due date and time: as announced on Blackboard.

(b) This assignment will be graded.

(c) Work in teams, as we have discussed at the beginning of the class.

\vspace{0.5cm}

\textbf{Programming project.}
\medskip

The goal of the project is to write a Java program that takes as input a DFA $M$ and an input string $w$,  simulates $M$ on $w$, and outputs ACCEPT if $M$ accepts $w$, and REJECT if $M$ does not accept $w$.


You will assume that the alphabet is $\Sigma = \{0,1\}$ and that the DFA has at most $20$ states. The states are labeled $1,2, \ldots, 20$.  The starting state is $1$.

The DFA is read from a file that has the following structure:

\begin{itemize}
  \item The first line has a single number $n$ which is the number of states.
  \item  The second line contains a list of numbers that indicate the accepting states. For example, it may be something like: 3 11 17 , meaning that the set of accepting states is $\{3, 11, 17\}$.
  \item The next lines indicate the transitions, with one transition per line. Each transition is the 3-tuple of numbers indicating (old state, symbol, new state).
        For example, it can be something like

        1  0 11

        1 1  7

        2 0 4

        2 1 6

        ...


        The 3-tuplet  1 0 11 says that from the state 1 with symbol 0 the machine moves to state 11,  the 3-tuplet  1 1 7 says that from the state 1 with symbol 1 the machine moves to state 7, and so on.

\end{itemize}

The program will ask the user to enter the input string $w$.

Test your program with the DFA  from Example 1.11, page 38 and the DFA from Example 1.68 (a), page 76, but  change the alphabet from $\{a,b\}$ to $\{0,1\}$.  Run  each  DFA on 3 strings that are accepted and 3 strings that are not accepted.
\bigskip

There should be a single submission per team, so each team should designate a submitter.

Submit on Blackboard 2 files, the first being the main one, and the second file being auxiliary.  \begin{enumerate}

  \item The 1st file (main component of your submission) is a pdf file that contains:
        \begin{itemize}
          \item The names of the people in the team.
          \item A short description of your program (1/2 page to 1 page should be enough), including a presentation of the data structures that you have used to represent the DFA and the input string w (probably some types of arrays or linked lists).
          \item The Java code of your program (should be well documented). The code also appears in the 2-nd  file  (see below) but I need it in the pdf format as well, so that I can make comments and observations.
          \item The  2  DFAs, that is the contents of the files read by your program.

          \item 2- 3 printouts with the screens when you run the program.
        \end{itemize}
  \item The 2-nd  file (auxilairy part of your submission)  is the JAVA file so that I can run your program.
\end{enumerate}
\medskip


\if01
  Submit on Blackboard:
  \begin{itemize}
    \item The Java code of your program (should be well documented). Also insert the Java code in the .pdf file, so that I can make comments and observations.
    \item The files that give the two DFAs.
    \item A short description of your program (1/2-1 page should be enough), including a presentation of the data structures that you have used to represent the DFA and the input string w (probably some types of arrays or linked lists).
    \item 2- 3 printouts with the screens when you run the program.
  \end{itemize}
\fi

\end{document}
